% Nouveau EGA V (cote 25): archive cote 25, site page 150 (PDF page 151).
% typescript. Draft compared: none.
% AI-transcribed by Claude Opus 5.5 (claude-opus-5-5, adaptive thinking, effort high), Message Batches API, batch ega25b (msgbatch_01A7n5YuXubUPNK1zsepi8RG), 2026-11-13; not yet checked by eye.
\scanpage{25/150}
\mspage{11}
% NOTE: a single handwritten diagonal stroke runs across the whole page, from the top right (above "conditions") to the bottom left (near "k >=1"); it cancels the whole text of the page, which stays legible under it.
d'ailleurs de vérifier ces dernières conditions \xout{xx}\inserted{pour} les fibres réduites
en les différents idéaux premiers (ou même seulement maximaux) de A/I .
\uline{Proposition} 8. Pour que l'augmentation A$\rightarrow$A/I soit \inserted{de \unsure{type}} quadratique
ordinaire, il f et s que les conditions suivantes soient vérifiées:
1º) \xout{xxx2} \xout{I/Ixxxxxx} \inserted[typed]{Les $gr^{n}_{I}(A)$ sont} projectifsde type fini sur A/I .

2º) Si d\xout{x} désigne le rang de $gr^{1}_{I}(A) = I/I^{2}$ en un $\underline{p} \in$Spec(A/I),
alors le rang de $gr^{n}_{I}(A)$ en ce même point est égal (ou simplement
$\geqslant$ ) à $\binom{n+d}{d} - \binom{n-2+d}{d}$ .

3º) (1º et 2º impliquent déjà que $\underline{K}$ = Ker $\underline{\theta}^{\unsure{2}}$ est un Module
inversible localement facteur direct dans \xout{xxxx} Sym$^{2}$(I/I$^{2}$) ). Le sous
Module $\underline{K}$ = Ker $\theta^{2}$ de Sym$^{2}$(I/I$^{2}$) est non dégénéré.

La démonstration est immédiate à partir du lemme suivant; \deleted{que je te
laisse le soin de démontrer:}

\xout{Propxsixxxx} \uline{Lemme} 8.1. Soient B un anneau, M un B-module projectif de
type fini, $\phi \in$ Sym$^{2}$(M) un élément ``non dégénéré'' , alors la multiplica-%
tion par $\phi$
\[ \mathrm{Sym}^{n\hand{\unsure{-2}}}(M) \longrightarrow \mathrm{Sym}^{n\deleted{\unsure{+2}}}(M) \]
est un isomorphisme sur un sous-module facteur direct, donc le \unsure{conoyau}
est un Module projectif de rang $\binom{n+d}{d} - \binom{n-2+d}{d}$ si M est de rang d.

On est immédiatement ramené au cas où B est le spectre d'un corps,
où c'est trivial (en fait, le ``non dégénéré'' nous a servi uniquement
\margin{||} pour assurer que les $\phi(x)$, $x \in$ Spec(B), sont tous $\neq 0$, et on pourrait
formuler le lemme plus raisonnablement en remplaçant 2 par un entier
k $\geqslant$1 quelconque, sans parler de non dégénérescence).

\uline{Proposition} 9. Soit A un anneau augmenté vers A/I, à augmentation
régulière, I/I$^{2}$ étant projectif de rang n sur A/I . Soit E un A-module

\mspage{150}
